% !TEX program = xelatex
%----------------------- Преамбула -----------------------
\documentclass[utf8x, 12pt, oneside, a4paper]{article}

\usepackage{extsizes} % Для добавления в параметры класса документа 14pt

% Для работы с несколькими языками и шрифтом Times New Roman по-умолчанию
\usepackage[english,russian]{babel}
\usepackage{fontspec}
\setmainfont{Times New Roman}

% ГОСТовские настройки для полей и абзацев
\usepackage[left=25mm,right=10mm,top=15mm,bottom=15mm]{geometry}
\usepackage{misccorr}
\usepackage{indentfirst}
\usepackage{enumitem}

\setlength{\parindent}{1.25cm}
\linespread{1.3}
\setlist{nolistsep} % Отсутствие отступов между элементами \enumerate и \itemize

% Дополнительное окружения для подписей
\usepackage{array}
\newenvironment{signstabular}[1][1]{
	\renewcommand*{\arraystretch}{#1}
	\tabular
}{
	\endtabular
}

% Переопределение стандартных \section, \subsection, \subsubsection по ГОСТу;
% Переопределение их отступов до и после для 1.5 интервала во всем документе
\usepackage{titlesec}

\titleformat{\section}[block]
{\bfseries\normalsize\filcenter}{\thesection}{1em}{}

\titleformat{\subsection}[hang]
{\bfseries\normalsize}{\thesubsection}{1em}{}
\titlespacing\subsection{\parindent}{\parskip}{\parskip}

\titleformat{\subsubsection}[hang]
{\bfseries\normalsize}{\thesubsubsection}{1em}{}
\titlespacing\subsubsection{\parindent}{\parskip}{\parskip}

% Работа с изображениями и таблицами; переопределение названий по ГОСТу
\usepackage{caption}
\captionsetup[figure]{name={Рисунок},labelsep=endash}
\captionsetup[table]{singlelinecheck=false, labelsep=endash}

\usepackage{graphicx}
\usepackage{slashbox} % Диагональное разделение первой ячейки в таблицах

% Цвета для гиперссылок и листингов
\usepackage{color}

% Гиперссылки \toc с кликабельностью
\usepackage{hyperref}

\hypersetup{
	linktoc=all,
	linkcolor=black,
	colorlinks=true,
}

% Листинги
\setsansfont{Arial}
\setmonofont{Courier New}

\usepackage{color} % Цвета для гиперссылок и листингов
\definecolor{comment}{rgb}{0,0.5,0}
\definecolor{plain}{rgb}{0.2,0.2,0.2}
\definecolor{string}{rgb}{0.91,0.45,0.32}

\usepackage{listings}
\lstset{
	basicstyle=\footnotesize\ttfamily,
	language=Python, % Или другой ваш язык -- см. документацию пакета
	commentstyle=\color{comment},
	numbers=left,
	numberstyle=\tiny\color{plain},
	numbersep=5pt,
	tabsize=4,
	extendedchars=\true,
	breaklines=true,
	keywordstyle=\color{blue},
	frame=b,
	stringstyle=\ttfamily\color{string}\ttfamily,
	showspaces=false,
	showtabs=false,
	xleftmargin=17pt,
	framexleftmargin=17pt,
	framexrightmargin=5pt,
	framexbottommargin=4pt,
	showstringspaces=false,
	inputencoding=utf8x,
	keepspaces=true
}

\DeclareCaptionLabelSeparator{line}{\ --\ }
\DeclareCaptionFont{white}{\color{white}}
\DeclareCaptionFormat{listing}{\colorbox[cmyk]{0.43,0.35,0.35,0.01}{\parbox{\textwidth}{\hspace{15pt}#1#2#3}}}
\captionsetup[lstlisting]{
	format=listing,
	labelfont=white,
	textfont=white,
	singlelinecheck=false,
	margin=0pt,
	font={bf,footnotesize},
	labelsep=line
}


% Годные пакеты для обычных действий
\usepackage{ulem} % Нормальное нижнее подчеркивание
\usepackage{hhline} % Двойная горизонтальная линия в таблицах
\usepackage[figure,table]{totalcount} % Подсчет изображений, таблиц
\usepackage{rotating} % Поворот изображения вместе с названием
\usepackage{lastpage} % Для подсчета числа страниц
\newcounter{numberlisting}
\setcounter{numberlisting}{1}
% ---------------------- Документ ----------------------
\begin{document}
	\begin{titlepage}
		\noindent\begin{minipage}{0.05\textwidth}
			\includegraphics[scale=0.3]{1.png}
		\end{minipage}
		\hfill
		\begin{minipage}{0.85\textwidth}\raggedleft
			\begin{center}
				\fontsize{12pt}{0.5\baselineskip}\selectfont 
				\textbf{Министерство науки и высшего образования Российской Федерации
				\\ Федеральное государственное бюджетное образовательное учреждение
				\\ высшего образования
				\\ <<Московский государственный технический университет
				\\ имени Н.Э. Баумана
				\\ (национальный исследовательский университет)>>
				\\ (МГТУ им. Н.Э. Баумана)}
			\end{center}
		\end{minipage}

		\begin{center}
			\fontsize{12pt}{0.1\baselineskip}\selectfont
			\noindent\makebox[\linewidth]{\rule{\textwidth}{4pt}} \makebox[\linewidth]{\rule{\textwidth}{1pt}}
		\end{center}
		\begin{flushleft}
			\fontsize{12pt}{1\baselineskip}\selectfont 
			ФАКУЛЬТЕТ \textbf{ИНФОРМАТИКА И СИСТЕМЫ УПРАВЛЕНИЯ}
			\\ КАФЕДРА \textbf{КОМПЬЮТЕРНЫЕ СИСТЕМЫ И СЕТИ (ИУ6)}
			\\ НАПРАВЛЕНИЕ ПОДГОТОВКИ \textbf{09.XX.XX \ Название направления}
			% бакалавры групп 1-3 и 6 – 09.03.01 Информатика и вычислительная техника;
			% бакалавры групп 4 и 5 – 09.03.03 Прикладная информатика;
			% магистры все – 09.04.01 Информатика и вычислительная техника
		\end{flushleft}
		\vfill
		\begin{center}
			\fontsize{16pt}{\baselineskip}\selectfont
			\textbf{О Т Ч Ё Т}
			\\ \fontsize{14pt}{\baselineskip}\selectfont \textbf{по лабораторной работе № <XX>}
		\end{center}
		\begin{table}[h!]
			\fontsize{14pt}{0.7\baselineskip}\selectfont
			\begin{signstabular}{p{3cm} p{13.4cm}}
				\textbf{Название:} & \uline{\fontsize{16pt}{\baselineskip}\selectfont{<Тема лабораторной работы>}\hfill}
				% Название по учебному плану: Практикум по программированию, 
				% Технологическая практика, Производственная практики, Преддипломная практика и т.п.
				\\ \ & \
				\\ \textbf{Дисциплина:} & \uline{\fontsize{16pt}{\baselineskip}\selectfont{<Название дисциплины без сокращений>}\hfill}
				% для учебной практики НУК ИУ МГТУ им. Н.Э. Баумана
			\end{signstabular}
		\end{table}
		\vfill
		\begin{table}[h!]
			\fontsize{14pt}{0.7\baselineskip}\selectfont
			\centering
			\begin{signstabular}[0.7]{p{3.25cm} >  {\centering\arraybackslash}p{4cm} > {\centering\arraybackslash}p{4cm} > {\centering\arraybackslash}p{4cm}}
				Студент & \uline{\hspace*{4cm}} & \uline{\hspace*{4cm}} & \uline{\hfill \textbf{<И.О. Фамилия>} \hfill} 
				\\ & \scriptsize (Группа) & \scriptsize (Подпись, дата) & \scriptsize (И.О. Фамилия)
				%дата обязательна, формат даты хх.хх.20хх
			\end{signstabular}

			\vspace{\baselineskip}

			\begin{signstabular}[0.7]{p{3.25cm} >  {\centering\arraybackslash}p{4cm} > {\centering\arraybackslash}p{4cm} > {\centering\arraybackslash}p{4cm}}
				Преподаватель & \ & \uline{\hspace*{4cm}} & \uline{\hfill \textbf{<И.О. Фамилия>} \hfill} 
				\\ & \ & \scriptsize (Подпись, дата) & \scriptsize (И.О. Фамилия)
				%дата обязательна, формат даты хх.хх.20хх
			\end{signstabular}
		\end{table}

		
		\begin{center}
			\normalsize Москва, <Год>
		\end{center}
	\end{titlepage}

	

	\normalsize
	\setcounter{page}{2}

	\addcontentsline{toc}{section}{ВВЕДЕНИЕ}
	ТЕКСТ ВВЕДЕНИЯ
	\\ ТЕКСТ

	ПРИМЕР НУМЕРОВАННОГО СПИСКА:
	\begin{enumerate} 
  		\item первый элемент первого уровня содержит список 
		\item  второй элемент первого уровня
	\end{enumerate}
	Пример ссылки на раздел: Подробная информация в разделе \ref{section1}
	\\ ПРИМЕР МАРКИРОВАННОГО СПИСКА:
	\begin{itemize}
		\item Each list item is ...
		\item Lists can be nested ...
			\begin{itemize}
				\item The maximum ...
				\item Switching ...
			\end{itemize}
		\item And so on.
	\end{itemize}
	
	Листинг \thenumberlisting \ - Код программы
	\begin{lstlisting}
		c = x^2
		print('test listing')
	\end{lstlisting}

	\begin{tabular}{ | p{100pt} | r | c | }
		\hline
		ИСЗ & Дата запуска & Масса, кг  \\ \hline
		Спутник-1 & 4 октября 1957 & 83.6 \\
		Спутник-2 & 3 ноября 1957 & 508.3 \\
		Эксплорер-1 & 1 февраля 1958 & 21.5 \\
		\hline
	\end{tabular}

	\begin{figure}[h]
		\centering
		\includegraphics[width=0.5\linewidth]{1.png}
		\caption{Тестовый рисунок}
		\label{fig:mpr}
	\end{figure}
	
	Формула: $c^{2}=a^{2}+b^{2}$
	\begin{math}
		c^2 = a^2 + b_1^2
	\end{math}
		\pagebreak
\end{document}
